\documentclass[12pt]{article}

\usepackage{fontspec}
\newfontfamily\handwriting{a-casual-handwritten-pen}[Path=../assets/fonts/,Extension=.ttf,Scale=1.5]
\newfontfamily\comic{Comic Neue}

\usepackage[utf8]{inputenc}
\usepackage{amsmath}
\usepackage{amssymb}
\usepackage{enumitem}
\usepackage{xcolor}
\usepackage{soul}
\usepackage{tikz}
\usepackage[most]{tcolorbox}
\usepackage{titlesec}
\usepackage{multicol}
\usepackage{environ}

\usetikzlibrary{fit, positioning, arrows.meta, decorations.pathreplacing}

% --- Custom Colors ---
\definecolor{primaryblue}{RGB}{42, 111, 151}
\definecolor{exbg}{RGB}{255, 240, 245}
\definecolor{rulebg}{RGB}{232, 245, 233}
\definecolor{highlight}{RGB}{255, 243, 176}
\definecolor{pinkanno}{RGB}{255, 105, 180}
\definecolor{purpleanno}{RGB}{147, 112, 219}
\definecolor{solutionbg}{RGB}{245, 245, 255}

% --- Header Styling ---
\titleformat{\section}{\color{primaryblue}\normalfont\Large\bfseries}{\thesection}{1em}{}
\titleformat{\subsection}{\color{primaryblue}\normalfont\large\bfseries}{\thesubsection}{1em}{}

% --- Friendly Boxes ---
% --- Auto-numbered Example Box ---
\newcounter{examplenum}

\newtcolorbox{friendlyexbox}[1]{
    colback=exbg,
    colframe=pinkanno!50,
    arc=5pt,
    title=\textbf{#1},
    coltitle=black,
    attach title to upper,
    after title={:\quad}
}

\newenvironment{friendlyex}{%
    \refstepcounter{examplenum}%
    \begin{friendlyexbox}{Example \theexamplenum}%
}{%
    \end{friendlyexbox}%
}

\newtcolorbox{friendlyrule}[1]{
    colback=rulebg,
    colframe=primaryblue!30,
    arc=5pt,
    fonttitle=\bfseries,
    title={#1},
    coltitle=primaryblue
}

\newcommand{\funhl}[1]{\sethlcolor{highlight}\hl{#1}}

% --- Show/Hide Solutions ---
\newif\ifshowsolutions

% Uncomment the next line to show solutions.
\showsolutionstrue

\newcommand{\solutionblank}{\vspace{25ex}}

\NewEnviron{solutionbox}{
\ifshowsolutions
\begin{tcolorbox}[
    colback=solutionbg,
    colframe=purpleanno!45,
    arc=5pt,
    title={Solution},
    coltitle=primaryblue,
    fonttitle=\bfseries
]
\BODY
\end{tcolorbox}
\else
\solutionblank
\fi
}

\begin{document}
\handwriting

\begin{center}
    \Large\textbf{\color{primaryblue}Module 4 Lesson 4\\Exponential Equations and Applications} \\
    \vspace{0.2cm}
    \hrule
\end{center}

\section*{Objectives}

\begin{friendlyrule}{In this lesson, we will learn how to}
\begin{itemize}
    \item solve exponential equations by matching bases;
    \item solve exponential equations using logarithms;
    \item apply exponential equations to compound interest problems.
\end{itemize}
\end{friendlyrule}

\section*{1. Equivalence Property}

\begin{friendlyrule}{Equivalence Property of Exponential Equations}
Let $a,x,y$ be real numbers with $a>0$ and $a\neq1$. Then
\[
a^x=a^y \qquad \text{implies that} \qquad x=y.
\]
\end{friendlyrule}

\textbf{Example:} \quad $5^x=25 \implies 5^x=5^2 \implies x=2$.

\begin{friendlyex}{}
Solve the following exponential equations.
\begin{enumerate}[label=(\alph*)]
    \item $3^x=27$
    \item $4^{2x-3}=32$
    \item $8^{x+2}=16^{x+1}$
\end{enumerate}
\end{friendlyex}

\begin{solutionbox}
\textbf{(a)} $3^x=27=3^3 \implies \boxed{x=3}$.

\textbf{(b)} Write both sides with base $2$: $4=2^2$, $32=2^5$.
\[
2^{2(2x-3)}=2^5 \implies 4x-6=5 \implies 4x=11 \implies \boxed{x=\frac{11}{4}}
\]

\textbf{(c)} Write both sides with base $2$: $8=2^3$, $16=2^4$.
\[
2^{3(x+2)}=2^{4(x+1)} \implies 3x+6=4x+4 \implies \boxed{x=2}
\]
\end{solutionbox}

\section*{2. Solving Exponential Equations Using Logarithms}

Some problems cannot be solved by setting the powers equal to each other (since the two sides cannot be written with matching bases). These are solved by applying a logarithm to both sides of the equation.

\begin{friendlyex}{}
Solve the following equations.
\begin{enumerate}[label=(\alph*)]
    \item $e^x=1200$
    \item $5^{2x-3}=17$
    \item $319=11^x+5$
\end{enumerate}
\end{friendlyex}

\begin{solutionbox}
\textbf{(a)} Apply $\ln$ to both sides:
\[
\ln(e^x)=\ln(1200) \implies x=\ln(1200)\approx \boxed{7.09}
\]

\textbf{(b)} Apply $\ln$ to both sides:
\[
(2x-3)\ln5=\ln17 \implies 2x-3=\frac{\ln17}{\ln5}\approx1.7604
\]
\[
2x\approx4.7604 \implies \boxed{x\approx2.38}
\]

\textbf{(c)} Isolate the exponential first:
\[
319-5=11^x \implies 314=11^x \implies x=\frac{\ln314}{\ln11}\approx \boxed{2.40}
\]
\end{solutionbox}

\begin{friendlyex}{}
Solve the following equations.
\begin{enumerate}[label=(\alph*)]
    \item $60=80e^{-0.4x}$
    \item $3^{x-2}=5^{x+1}$
\end{enumerate}
\end{friendlyex}

\begin{solutionbox}
\textbf{(a)} Isolate the exponential, then apply $\ln$:
\[
\frac{60}{80}=e^{-0.4x} \implies 0.75=e^{-0.4x} \implies \ln(0.75)=-0.4x
\]
\[
x=\frac{\ln(0.75)}{-0.4}\approx \boxed{0.72}
\]

\textbf{(b)} Apply $\ln$ to both sides:
\[
(x-2)\ln3=(x+1)\ln5
\]
\[
x\ln3-2\ln3=x\ln5+\ln5
\]
\[
x(\ln3-\ln5)=\ln5+2\ln3
\]
\[
x=\frac{\ln5+2\ln3}{\ln3-\ln5}\approx\frac{1.6094+2.1972}{1.0986-1.6094}=\frac{3.8066}{-0.5108}\approx \boxed{-7.45}
\]
\end{solutionbox}

\section*{3. Application: Compound Interest}

\begin{friendlyex}{}
If \$5000 is put into an interest-bearing account at an interest rate of $2.2\%$ compounded monthly, find the time required for the account to \textbf{earn} \$1000. Round to the nearest month.
\end{friendlyex}

\begin{solutionbox}
Earning \$1000 means the account grows from \$5000 to \$6000, so $A=6000$, $P=5000$, $r=0.022$, $n=12$.
\[
6000=5000\left(1+\frac{0.022}{12}\right)^{12t}
\]
\[
1.2=\left(1.0018333\right)^{12t}
\]

Apply $\ln$ to both sides:
\[
\ln(1.2)=12t\,\ln(1.0018333)
\]
\[
12t=\frac{\ln(1.2)}{\ln(1.0018333)}\approx\frac{0.18232}{0.0018317}\approx99.55
\]

Rounding to the nearest month, $12t\approx \boxed{100 \text{ months}}$ (about $8$ years and $4$ months).
\end{solutionbox}

\section*{Summary}

\begin{friendlyrule}{Key Ideas}
\begin{itemize}
    \item Equivalence Property: if $a^x=a^y$ (same base), then $x=y$. Rewrite both sides with a common base whenever possible.
    \item When the two sides cannot share a common base, apply $\ln$ (or $\log$) to \textbf{both sides} of the equation, then use the power rule $\ln(b^x)=x\ln(b)$ to bring the variable down out of the exponent.
    \item Always \textbf{isolate} the exponential expression first before applying a logarithm (e.g.\ subtract or divide off any extra terms/factors).
    \item When variable exponents appear on both sides with different bases (like $3^{x-2}=5^{x+1}$), take $\ln$ of both sides, expand with the product rule, then collect all $x$-terms on one side to solve.
\end{itemize}
\end{friendlyrule}

\end{document}
